\documentclass[11pt]{article}
\usepackage{preprint}

\title{The Significance of Law and Agreements in International Relations}
\author{Jiadi He}
\shorttitle{Law and Agreements in International Relations}
\keywords{international relations; law; institutions}

\begin{document}
\maketitle
\keywordsline

Over the course of the semester, we covered critical topics on international relations. Each topic highlighted a gap in international study that is critical in advancing the discipline. The weekly critiques provided a scope for analyzing different perspective on traditional practices on international relations. The critiques highlighted the various ways in which the current study of international relations is limited, and provided an avenue to assess the potential opportunities at addressing these limitations for the benefit of future research. Over the past century, the relationship of actors in international relations has evolved leading to the gap in the discipline. These relationships influence key decisions that shape the landscape of international politics. These aspects signify why law and agreements was the most important topic during the semester. The topic covered readings that analyzed international law, global organizations and institutions, and negotiating and bargaining, which are critical elements in modern international relations. The topic on law and agreements is essential to international relations as it provides an opportunity to define the legal framework of state and interstate developments and to regulate the power exerted by international actors by protecting the interest of involved parties.

Law and agreements help in redefining the traditional outlook of the power occupied by the state in international relations. Lake (2007) highlighted those states as the central actors in international relations through state-centrism. His research was critiqued in the sense that it had a fixed approach in understanding international relations, which limited the scope of the study. The role of the state in international relations has evolved over the past century as more actors have an influence on key political decisions that shape international politics. According to Biersteker (2013), the concepts of states, sovereignty, and territory are no longer fixed and constant as depicted in traditional international relations. The three concepts continuously evolve depending on the region and timeline in history. This factor demystifies the fixed approach of considering states as central actors in international relations because it lacks a quick adaptation to changes in modern politics. Law and agreements provided a scope that addressed this challenge. For example, Martin and Simmons (2013) showed how international organizations provided states with opportunities to address critical issues such as high costs and information deficits. The two scholars also addressed how such actions, by international actors, influence state behaviors, which is critical in defining state relationships with other actors on an international landscape. Law and agreements are critical in expanding the scope of international study from a state-centered to a global perspective that incorporate key international relationships.

Law and agreements are an essential tool in defining the power that actors have in international relations. States are considered the central actors in international relations, which is supported by societal perceptions. The authority and power held by states is acquired through social constructs that evolve over a period of time. States exert their dominance through political and social organizations, and the niche they occupy on an international level. Biersteker (2013) highlights that other actor’s actions, practices, and regulations are responsible for defining and redefining the state, sovereignty, and territory. However, traditional international relations that use state-centric theories only primarily focus on a single political order that is characterized by domestic ideals and institutions. The international scope has significantly evolved, redefining the roles of critical international actors. International law serves as a distinguishing feature in international relations. The attributes of international law distinguish practices, political agreements, and institutions through its binding principles that are globally accepted and applicable (Simmons, 2013). For example, public international law provides a set of binding rules that regulate interactions between states. These laws can be used to redefine power and establish critical roles for states with consideration of their status quo in international relations.

International politics lacks universal ethical standards and principles, which gives states and other international actors freedom of exerting their power beyond domestic boundaries with limited restrictions. Powerful states leverage the power they have in terms of resources to gain access and control essential services in foreign states, which undermines the growth small states can accomplish when empowered. The regulation of power in international politics was undermined in traditional international relations, giving actors significant freedom to operate. Wohlforth (2007) highlights four assumptions of realism in traditional international relations that undermined other objectives of analyzing power. Traditional realism relies on these four assumptions, namely; Groupism, egoism, anarchy, and human inequalities in power which propose a primary focus on interstate relationships and the balancing of power. However, under realism, states utilize policies that are different according to the region that the state exercises its power (Donnelly, 2000). This factor can be attributed to the exercises taking up a practical stance instead of a moral one. These issues highlight a critical gap in addressing the challenges faced by international actors with limited international power. International law and agreements provide an avenue of bridging the gap and protecting the interests of each international actors. The tides of the Cold War paved way for the incorporation of transnational actors in international relations, which compliments the influence international law will have in addressing the issues (Risse, 2013). Regulating international power is critical in maintaining peace, law, and order that was undermined during the persistent war periods in the past century.

International law is critical in regulating power of international actors through accepted principles that are also designed to mediate conflict between the actors and avoid prolonged disputes that lead to war (Simmons, 2013). Notably, the elements of war were a regular occurrence in the past century, which required a sound approach to diffuse the global tension that rose during the aftermath. Conquest in war was considered a primary tactic in advancing international interest especially in the 20\textsuperscript{th} century. Powerful states wanted to strategically place themselves in the international political landscape, which required a big statement to cement the position. Levy (2013) highlights that majority of realist pointed out that war resulted from predatory aggression through states that believed military force was the logical way of advancing their interests. Conflict spirals between dissatisfied states either inadvertently or intentionally aimed at maintaining the security of the state was also mentioned as a key contributor to war. Most states engaged in war due to the uncertainty of what the rival state may choose to do in the event of a prolonged dispute. While each state reserves the right to protect its interest and citizens, war led to massive losses in the economy and deterioration in the society (Simpson, 2008). The effects of war were difficult to reverse and left many countries at a disadvantage that still influences their political power play in modern international relations (Moravcsik, 2008). Thus, international law is critical in addressing such issues. The laws incorporate punishments for states that abuse their international powers to undermine smaller states which is critical in safeguarding international interests. Law and agreements are a basis of resolving international conflict that was challenging in international relations.

Laws and agreements are critical in determining the impact of diffusion process on state policies. The diffusion process involves both domestic and international actors who rely on their power in making critical decisions. The relationship between states and other international actors are responsible for the critical decisions in international politics that are shaped through mutual benefits and conflicts. The conflicts and mutual benefits arise from the international dependence exhibited by majority of the nations around the world (Zurn, 2013). The interaction between states depend on how compatible their political and social systems are in order to compliment the relationship and sustain collaborations. The collaborations require a high level of diffusion which occurs at the international and domestic level (Gilardi, 2013). The national government is responsible for diffusing policies that determined how the nation is ran and regulates domestic interactions. The traditional approach in international relations depicts that the national government is solely responsible for making decisions that shape domestic integration. However, transnational diffusion influence on these decisions has increased as the role of international actors evolve. Domestic decisions made in one country influence how various foreign policies, ideals and norms are shaped (Gilardi, 2013). This factor is attributed on the influence private international actors have as secondary decision markers. Secondary decision makers, with significant political power, dictate domestic and international policies in favor of their interests which undermines the roles of national governments and other actors in international politics. The diffusion process threads policies in international politics which are difficult to reverse especially if the socially constructed notions of power are deeply rooted. International law highlights the role of private actors such as international organisations and institutions, and provides an avenue to scrutinize their influence on the dependence between states and other actors, safeguarding individual and collective interests (Simmons, 2013).

Law and agreements provide insight on how countries operate in free economies. Free economies allow countries to operate with minimal restrictions to promote economic integration and growth. Successful diffusion and integration between states and other actors lead to new collaborations that's are unchartered in international relations. Trade activities between states have significantly diversified through the diffusion of policies which have opened up more countries to flows of exports and imports in large volumes (Milner, 2013). Technological advancements and diffusion of government policies have increased the collaboration between developing and developed states in the race towards economic dominance. Milner (2013) highlights the significance of successful negotiations such as the GATT, a multilateral trade negotiation of 1994. The negotiation influenced the rush towards free trade aimed at establishing a free world economy that benefits involved states. However, the risk in a free economy is the consistently shifting balance of power among states. Certain states such as China are slowly approaching “developed country” status, which means their influence will increase, creating a need to renegotiate their terms in international trade. States leverage their benefits in international relationships through negotiations to maximize on their roles and achieve better terms in international collaborations (Odell, 2013). Simmons (2013) highlights how international law is developed through conceptualization of legal, judicial, constitutional, and global administrative law that are also essential in defining trade practices. Understanding the free economy requires a legal frame that can be developed through law and agreements shaped to regulate international trade.

The legal framework on trade trickles down to the financial aspects in international relations, which determine how money is exchanged between actors. States trade according to relationship and regulations stipulated within the operating market. Governments and international organization have a significant influence on foreign exchange rates, financial regulations, and financial aid, which dictates whether international actors achieve economic destitution or prosperity (Tomz, 2013). Financial activities of international firms and foreign immigrants contribute to the financial ability of an international actor, which highlights the significance of understanding international finance in the global economy. Financial laws vary across regions and over time as countries try to safeguard their investments while also benefiting from international financial activities. Countries that wish to improve their domestic economy may limit foreign investment through harsh taxes, while some countries may set policies that allow free flow of investment between borders. Traditional international relations gave primary focus to financial considerations made during the world wars which included capital barriers and leveraging on acquiring equity through foreign investments (Tomz,2013). However, most capital barriers between states have been removed leading to free flow of investments in most regions. International organizations also provide financial aid to needy states, aimed at improving their economic situations and safeguard interstate collaboration (Martin \& Simmons, 2013). These factors highlight the need of understanding international law and agreements that define the relationship between actors and the financial implications they have on states and international institutions.

International laws and agreements provide a scope for analyzing the role of global institutions. The role of global institutions is scoped into the laws and agreements that govern their operations. Global institutions are tasked with the responsibility of delivering the collective goals of member states and provide support in critical situations. However, their jurisdiction has been put into questions by realists who believe they mostly deal with low political domain issues (Stein, 2008). This issue is problematic because most states deal with high political issues that threaten global security and collaborations. It calls into question whether global institutions give promises on dealing with high political issues that cannot be delivered. These faults are heightened by the rights given to member states in the institutions. States within a global institution should share equal rights as the collective goals and objectives are meant to deliver mutual benefit to each member. However, the level of power of states in the global spectrum influences how each state is treated in global institutions. In some instances, powerful states dictate the running of global institutions to serve their will and undermine the position of other member states (Richardson, 2008). It is critical to understand the jurisdiction of each institution and how they shape the global landscape with regards to the relationship between states (Martin \& Simmons, 2013).

The most critical aspect served by the topic is highlighting the significant changes to the international legal system that had not been considered in traditional international relations. Byers (2010) highlights that the changes in the international legal system are complimented by traditional rules, judicial models, treaties, acceptable principle, and scholarly items. Politicians are the primary domestic law makers, while diplomats are internationally recognized lawmakers, complimented by international organizations and corporations (Byers, 2010). These factors update the traditional international relations notion on states being the primary actors. It also compliments transnational studies which have been which neglected in traditional international relations. The gap between international relations and modern concepts have furthered global issues such as environmental problems that are disregarded as non-critical in international relations (Mitchell, 2013). Governments and international actors expected to deal with these issues through the collaborations created through global interaction. The changes in the international legal system provide an opportunity to understand the positioning of each actor and define the roles they are expected to play in addressing critical issues. Changes in the relationship between international actors also influence how the international legal system evolves over time depending on the region or jurisdiction of the global institution.

Analyzing the course reading for the semester provides a wide scope of understanding the relationship between the topics and the developments in international relations. Traditional international relations have a significant gap to cover and incorporate the values and principles of modern international relations. However, the critical aspect that stood out in the analysis is the influence international law and agreements have in international relations. The relationship between states and international actors is consistently evolving and leading to unchartered territory, which needs to be covered through a framework that protects the value of international politics. The centrality of states in international relations has been demystified which provides an opportunity to curve out new niches for international actors, aiming to achieve breakthrough in international study. A proper understanding of the current international law and agreements provides scholars with an avenue of expanding international studies, and complimenting the collaborations between international actors into fruitful relationships that respect the power of each actor.

\begin{references}
\item Biersteker, T. State, Sovereignty, and Territory (2013). \textit{Handbook of international relations} (2nd ed., pp. 245–272). essay, SAGE.
\item Byers, M. (2010). International Law. In \textit{The Oxford Handbook of International Relations} (pp. 612-634). Oxford University Press.
\item Donnelly, J. (2000). \textit{The Oxford Handbook of international relations} (pp. 150-162). essay, Oxford University Press
\item Gilardi, F. (2013). Transnational Diffusion: Norms, Ideas and Policies. In \textit{Handbook of international relations} (pp. 453-477). SAGE.
\item Lake, D. (2007). The State and International Relations. In \textit{The Oxford Handbook of international relations} (pp. 41–61). essay, Oxford University Press.
\item Levy, J. (2013). Interstate War and Peace. In \textit{Handbook of international relations} (pp. 581-606). SAGE.
\item Martin, L., \& Simmons, B. (2013). International Organizations and Institutions. In \textit{Handbook of international relations} (pp. 326-351). SAGE.
\item Milner, H. (2013). International Trade. In \textit{Handbook of international relations} (pp. 720-745). SAGE.
\item Mitchell, R. (2013). International Environmental Politics. In \textit{Handbook of international relations} (pp. 801-826). SAGE.
\item Moravcsik, A. (2008). The New Liberalism. In The \textit{Oxford Handbook of international relations} (pp. 234-254). essay, Oxford University Press.
\item Odell, J. (2013). Negotiation and Bargaining. In \textit{Handbook of international relations} (pp. 379-400). SAGE.
\item Richardson, J., B. (2008). The Ethics of Neoliberal Institutionalism. In \textit{The Oxford Handbook of international relations} (pp. 222-233). essay, Oxford University Press.
\item Risse, T. (2013). Transnational Actors and World Politics. In \textit{Handbook of international relations} (pp. 426-477). SAGE.
\item Simmons, B. (2013). International Law. In \textit{Handbook of international relations} (pp. 352-378). SAGE.
\item Stein, A. (2008). Neoliberal Institutionalism. In \textit{The Oxford Handbook of international relations} (pp. 201-221). essay, Oxford University Press.
\item Tomz, M. (2013). International Finance. In \textit{Handbook of international relations} (pp. 692-719). SAGE.
\item Wohlforth, W. (2007). Realism. In \textit{The Oxford Handbook of international relations} (pp. 131-149). essay, Oxford University Press.
\item Zurn, M. (2013). Globalization and Global Governance. In \textit{Handbook of international relations} (pp. 401-425). SAGE.
\end{references}

\end{document}
